% ---
% title: Lost cards: kings
% tags: generalized-bayes, total-probability, complement
% ---

\begin{problem}
	Two cards have been lost from a deck of 52 cards; we do not know which ones.
	We draw 4 of the remaining 50 cards at random.
	(a) Find the probability that at least one king is drawn.
	(b) If at least one king is drawn, what is the probability that both lost
	cards were kings?
\end{problem}

\begin{solution}
	\idea We do not know how many kings were lost, so we split the situation into
	three cases by the number of lost kings, as in
	Exercise~\exref{lost-cards-aces}. In each case ``at least one king'' is
	easiest through its complement, ``no king''. Then (a) is the law of total
	probability over the three cases, and (b) turns the conditioning around with
	the generalized Bayes formula.

	\medskip
	\textbf{The three cases.} Let $B_0$, $B_1$, $B_2$ be the events that $0$,
	$1$ or $2$ of the lost cards are kings. Exactly one of them happens, so they
	are mutually exclusive and their union is $\Om$. We know nothing about the
	lost cards, so every pair of the 52 cards is equally likely to be the lost
	pair. There are $\binom{52}{2}=1326$ pairs. A pair without a king is chosen
	from the 48 non-kings, a pair with one king is one of the 4 kings and one of
	the 48 non-kings, and a pair of kings is chosen from the 4 kings:
	\[ P(B_0)=\frac{\binom{48}{2}}{1326}=\frac{1128}{1326}=\frac{188}{221},\quad
		P(B_1)=\frac{4\cdot48}{1326}=\frac{192}{1326}=\frac{32}{221},\quad
		P(B_2)=\frac{\binom{4}{2}}{1326}=\frac{6}{1326}=\frac{1}{221}. \]
	Check: $188+32+1=221$, so the three probabilities add up to $1$.

	\medskip
	\textbf{The draw in each case.} Let $A$ be the event that at least one king
	is drawn. Its complement $A^c$ is ``no king is drawn''. If $i$ kings were
	lost, the remaining 50 cards contain $4-i$ kings and $46+i$ non-kings. The
	draw of 4 cards is one of $\binom{50}{4}=230300$ equally likely sets, and
	``no king'' means that all 4 cards come from the $46+i$ non-kings:
	\[ P(A^c\mid B_i)=\frac{\binom{46+i}{4}}{\binom{50}{4}}. \]
	So
	\[ P(A^c\mid B_0)=\frac{163185}{230300}\approx0.7086,\quad
		P(A^c\mid B_1)=\frac{178365}{230300}\approx0.7745,\quad
		P(A^c\mid B_2)=\frac{194580}{230300}\approx0.8449. \]
	Given $B_i$, the conditional probabilities of $A$ and $A^c$ add up to $1$,
	so $P(A\mid B_i)=1-P(A^c\mid B_i)$. In exact form:
	\[ P(A\mid B_0)=\frac{13423}{46060}\approx0.2914,\quad
		P(A\mid B_1)=\frac{221}{980}\approx0.2255,\quad
		P(A\mid B_2)=\frac{38}{245}\approx0.1551. \]
	The more kings were lost, the fewer kings are left, and the smaller the
	chance of drawing one.

	\step{(a)} The event $A$ is the union of the three disjoint events $AB_0$,
	$AB_1$, $AB_2$, so their probabilities add up. By the multiplication rule
	$P(AB_i)=P(B_i)P(A\mid B_i)$:
	\[ P(AB_0)=\frac{188}{221}\cdot\frac{13423}{46060}=\frac{13423}{54145}\approx0.2479, \]
	\[ P(AB_1)=\frac{32}{221}\cdot\frac{221}{980}=\frac{8}{245}=\frac{1768}{54145}\approx0.0327, \]
	\[ P(AB_2)=\frac{1}{221}\cdot\frac{38}{245}=\frac{38}{54145}\approx0.0007. \]
	Adding them (the law of total probability):
	\[ P(A)=\frac{13423+1768+38}{54145}=\frac{15229}{54145}\approx0.2813. \]

	\step{(b)} We need $P(B_2\mid A)$: among the draws with at least one king,
	the share of those where both lost cards were kings. By the definition of
	conditional probability, the numerator is $P(AB_2)$ and the denominator is
	$P(A)$, both found in (a):
	\[ P(B_2\mid A)=\frac{P(B_2)P(A\mid B_2)}{P(A)}
		=\frac{38/54145}{15229/54145}=\frac{38}{15229}\approx0.0025. \]

	\begin{answer}
		(a) $\dfrac{15229}{54145}\approx0.2813$ \quad
		(b) $\dfrac{38}{15229}\approx0.0025$
	\end{answer}
	\begin{remark}
		The answer to (a) is the same as with a full deck:
		\[ 1-\frac{\binom{48}{4}}{\binom{52}{4}}=\frac{15229}{54145}. \]
		This is no accident. We
		never look at the lost cards, so they are just two more cards taken at
		random, and the 4 drawn cards are still 4 random cards of the original
		deck.

		In (b) the posterior $0.0025$ is about half of the prior
		$P(B_2)=1/221\approx0.0045$. Drawing a king is a sign that kings are
		still in the deck, so it makes ``both lost cards were kings'' less
		likely.

		\textbf{Common mistake:} answering (b) with the prior
		$P(B_2)=1/221\approx0.0045$, or with $P(A\mid B_2)\approx0.1551$,
		instead of the posterior $P(B_2\mid A)$. The question fixes $A$ and asks
		about $B_2$, so $A$ goes after the bar.
	\end{remark}
\end{solution}
